\documentclass[conference]{IEEEtran}
\IEEEoverridecommandlockouts

\usepackage[utf8]{inputenc}
\usepackage[T1]{fontenc}
\usepackage{amsmath,amssymb}
\usepackage{graphicx}
\usepackage{booktabs}
\usepackage{cite}
\usepackage{xcolor}
\usepackage[hidelinks]{hyperref}

\begin{document}

\title{Protecting the Right Weights on a Noisy Hybrid SLC/MLC ReRAM Chip}

\author{
\IEEEauthorblockN{Author One}
\IEEEauthorblockA{\textit{Department} \\
\textit{Institution}\\
City, Country \\
email@example.com}
\and
\IEEEauthorblockN{Author Two}
\IEEEauthorblockA{\textit{Department} \\
\textit{Institution}\\
City, Country \\
email@example.com}
}

\maketitle

\begin{abstract}
Analog resistive-RAM (ReRAM) processing-in-memory computes matrix--vector products inside the memory
array and removes most of the data movement that dominates transformer inference. Its weakness is that
a cell cannot be programmed exactly. A single-level cell (SLC) stores one bit and tolerates this error,
while a multi-level cell (MLC) stores two bits and is denser but far less robust. A hybrid chip keeps
the important weights in SLC and the rest in MLC, so the central question is which weights to protect.
We study this question for DistilBERT and BERT-Base on SST-2 and MNLI with a cell-level crossbar simulator that
models 8-bit weights, differential pairs, bit-slicing, programming noise, stuck-at faults, read noise,
a fixed ADC and IR drop. After a singular value decomposition (SVD) of every weight matrix, we compare
six placements of the SVD components: all-SLC, all-MLC, random, largest singular value (magnitude),
a HyFlexPIM-style loss gradient, and a Fisher-information estimate of the damage each component would
suffer in MLC. Storing everything in MLC fails at realistic noise, and BERT-Base falls from about 90\% to
near chance. Protecting a small share of components recovers almost all of the loss, but only if the right
components are chosen. Ranking by singular value matches the best placement we could construct,
whereas the gradient placement loses 3--12 accuracy points at high noise. Stuck-at cells are the one
failure that placement does not fix: all-MLC fails at 0.1\% stuck cells and all-SLC at 1\%. Fault-aware
mapping, which reorders the components around a chip's known stuck cells, recovers much of this loss. The
same results hold on MNLI.
\end{abstract}

\begin{IEEEkeywords}
processing-in-memory, ReRAM, multi-level cell, transformer, singular value decomposition, fault tolerance
\end{IEEEkeywords}

\section{Introduction}
Transformer models \cite{devlin2019bert} are expensive to run, mainly because billions of weights must be
moved between memory and the processor. Processing-in-memory (PIM) removes this movement by computing where
the weights are stored. Analog ReRAM crossbars are an efficient form of PIM, but a cell cannot be
programmed exactly, so every stored weight carries error. Dense multi-level cells (MLC) suffer most from
this error, and reliable single-level cells (SLC) are costly. A hybrid chip therefore keeps only the most
important weights in SLC, and the open question is which weights these are.

This paper studies that question for SVD-factorised transformers on a simulated hybrid SLC/MLC ReRAM chip.
Section~II gives the background, including the HyFlexPIM method \cite{song2025hyflexpim} and our
contributions. The remaining sections will describe the simulator, the experiments and the results.

\section{Background}
\subsection{Analog ReRAM Crossbars}
A ReRAM crossbar computes a matrix--vector product inside memory. Each weight is stored as the
conductance of resistive cells. A voltage proportional to each input is applied to a row, and the currents
on each column add up to the output. This avoids most of the data movement of conventional hardware, and it
computes a whole column in one step.

The cost is accuracy. A cell cannot be programmed to an exact conductance, and cells can also be stuck at
a fixed conductance. A single-level cell (SLC) stores one bit in two well-separated levels and tolerates
considerable error. A multi-level cell (MLC) stores two bits in four closely spaced levels. It needs fewer
cells per weight, which gives higher density, throughput and energy efficiency, but the same absolute
conductance error is a larger error in the stored value. Transformers are sensitive to such error, and
running BERT entirely on two-bit cells can cost tens of accuracy points \cite{song2025hyflexpim}.

\subsection{HyFlexPIM}
HyFlexPIM \cite{song2025hyflexpim} is a hybrid SLC/MLC mixed-signal PIM accelerator. A single analog PIM
module can switch between SLC and MLC operation with under 1\% area and energy overhead, so critical
weights can be stored in SLC and the remaining weights in the denser MLC. Brute-force mapping on such a chip
gives little benefit, because too few weights are error-tolerant enough for MLC. HyFlexPIM therefore adds an
algorithmic step called \emph{gradient redistribution}:
\begin{enumerate}
\item Factorise each weight matrix with singular value decomposition (SVD) into components ordered by
singular value.
\item Truncate the smallest components to remove redundancy.
\item Fine-tune the truncated model, which concentrates the loss gradient in a small subset of components.
\item Rank the components by the absolute gradient of the loss with respect to their singular value.
\item Store the top $k\%$ in SLC and the remaining components in MLC.
\end{enumerate}
All of these steps run in software before deployment, and the hardware only receives the final truncated
matrices. HyFlexPIM models MLC error by multiplying each weight by $(1+\eta)$ with Gaussian $\eta$, with the
variance chosen to match a bit error rate of about 4.04\% measured on fabricated MLC RRAM. It reports that
protecting 5--30\% of components in SLC keeps most GLUE tasks within 1\% of the noise-free baseline, and that
gradient-based selection outperforms magnitude- and rank-based selection on MRPC and CoLA with BERT-Base.

\subsection{Open Questions and Our Contributions}
The evidence for gradient-based selection covers a small number of tasks and a single noise model, and the
model has no stuck-at cells. It is unclear whether the rule holds under other error models, whether a
cheaper rule performs as well, and how stuck-at cells affect each placement. Our contributions address
these points:
\begin{itemize}
\item \textbf{A crossbar simulator.} A cell-level PyTorch model of SVD-factorised transformers on SLC and
MLC crossbars, covering weight quantisation, differential pairs, bit-slicing, programming noise, stuck-at
faults, read noise, a fixed ADC and IR drop. All placements are evaluated on the same simulated chips.
\item \textbf{Criticality-aware component placement.} A comparison of six policies for choosing which SVD
components are stored in SLC: all-SLC, all-MLC, random, singular-value magnitude, the HyFlexPIM gradient,
and a Fisher-information estimate of the damage each component would suffer in MLC.
\item \textbf{Variation with SLC share.} How the accuracy of each placement policy changes as the share of
components stored in SLC increases.
\item \textbf{Fault analysis.} An evaluation of how programming noise and stuck-at cells, which HyFlexPIM
does not model, affect each placement policy.
\item \textbf{Fault-aware mapping.} A post-test remapping that uses the freedom to store SVD components in
any order: once a chip's stuck cells are known, each component is assigned to the physical slot where it
would raise the loss least, at no hardware cost.
\end{itemize}
We study DistilBERT and BERT-Base on the SST-2 task of GLUE, and repeat the main experiments on MNLI.

\section{Methodology}
Our pipeline has three stages: SVD factorisation of the transformer, a cell-level crossbar simulation, and
the placement of the SVD components on SLC and MLC cells. This section describes the first stage.

\subsection{SVD of the Transformer}
\subsubsection{Models and dataset}
We use two encoder transformers, DistilBERT \cite{sanh2019distilbert} and BERT-Base \cite{devlin2019bert},
and the SST-2 sentiment task of GLUE \cite{wang2019glue}. SST-2 is a binary classification task. It has
67,349 training sentences and 872 validation sentences, and all accuracies in this paper are measured on the
validation set. Both models start from public checkpoints that are already fine-tuned on SST-2
(\texttt{distilbert-base-uncased-}\allowbreak\texttt{finetuned-sst-2-english} and
\texttt{textattack/}\allowbreak\texttt{bert-base-uncased-SST-2}). Sentences are truncated to 128 tokens.

The second task is MNLI, a three-class entailment task on sentence pairs. It has 392,702 training pairs and
9,815 matched validation pairs. To keep the simulations affordable, the chip experiments on MNLI use a fixed
random subset of 2,000 validation pairs. The checkpoints are
\texttt{typeform/}\allowbreak\texttt{distilbert-base-uncased-mnli} and
\texttt{textattack/}\allowbreak\texttt{bert-base-uncased-MNLI}. The label order of each checkpoint is matched
to GLUE on the evaluation subset.

\subsubsection{Factorisation}
Following HyFlexPIM, we factorise every linear layer of the encoder blocks (the query, key, value and
output projections and the two feed-forward layers) as
\begin{equation}
W \approx U\,\mathrm{diag}(\sigma)\,V^{\top}, \qquad W\in\mathbb{R}^{m\times n},
\end{equation}
which gives 36 factorised layers for DistilBERT and 72 for BERT-Base. We keep the $k=\lfloor mn/(m+n)\rfloor$
largest singular values (for example $k=384$ for a $768\times768$ layer), the largest rank for which the two
factors $U$ and $\sigma V^{\top}$ together hold no more parameters than $W$. On the chip each layer then
becomes two crossbars, $x\rightarrow[\sigma V^{\top}]\rightarrow[U]\rightarrow y$, and component $r$ consists
of row $r$ of $\sigma V^{\top}$ together with column $r$ of $U$. This gives 16,584 components for DistilBERT
and 33,168 for BERT-Base.

\subsubsection{Recovery fine-tuning}
Truncating to rank $k$ removes information, so we fine-tune the factorised model on the SST-2 training set
to recover the lost accuracy. We use AdamW with a learning rate of $2\times10^{-5}$ decayed linearly to
zero, batch size 32, 8 epochs, mixed-precision (fp16) autocast and seed 0. On MNLI we use 2 epochs on a
random 100,000 of the training pairs. Table~\ref{tab:svd} shows the validation accuracy of the
SVD-factorised models after fine-tuning.

\begin{table}[t]
\caption{Validation accuracy (\%) of the original and the SVD-factorised models (full validation sets).}
\label{tab:svd}
\centering
\begin{tabular}{lcccc}
\toprule
 & \multicolumn{2}{c}{SST-2} & \multicolumn{2}{c}{MNLI} \\
\cmidrule(lr){2-3}\cmidrule(l){4-5}
Model & Original & SVD & Original & SVD \\
\midrule
DistilBERT & 91.06 & 89.11 & 82.21 & 80.76 \\
BERT-Base  & 92.43 & 90.14 & 84.57 & 82.87 \\
\bottomrule
\end{tabular}
\end{table}

All later experiments use these fine-tuned factorised models as their noise-free reference, so the
accuracies in the tables that follow should be read against 89.11\% (DistilBERT) and 90.14\% (BERT-Base)
on SST-2, and 80.9\% and 83.2\% on the 2,000-pair MNLI subset.

\subsection{INT8 Quantisation}
A crossbar stores integers, so every weight and every input is quantised to 8 bits before it reaches the
array. Weights use a symmetric 8-bit scale per output row, and inputs use an 8-bit scale per tensor. In the
simulator, an ideal crossbar with no noise or faults reproduces an INT8 matrix product exactly (relative
error below $10^{-4}$). The noise-free all-SLC chip therefore measures the accuracy of the INT8 model, and
every difference between a noisy chip and this value is caused by the cell errors alone.

Table~\ref{tab:int8} compares floating-point (FP32) and INT8 accuracy for the original checkpoints and for
the fine-tuned SVD models.

\begin{table}[t]
\caption{SST-2 validation accuracy (\%) in FP32 and INT8, with no hardware error.}
\label{tab:int8}
\centering
\begin{tabular}{lcccc}
\toprule
 & \multicolumn{2}{c}{Original} & \multicolumn{2}{c}{SVD + fine-tune} \\
\cmidrule(lr){2-3}\cmidrule(lr){4-5}
Model & FP32 & INT8 & FP32 & INT8 \\
\midrule
DistilBERT & 91.06 & 90.37 & 89.11 & 88.99 \\
BERT-Base  & 92.43 & 91.86 & 90.14 & 89.91 \\
\bottomrule
\end{tabular}
\end{table}

INT8 quantisation costs 0.6--0.7 points for the original models and 0.1--0.2 points for the SVD models.
It is therefore a small effect compared with the 1.9--2.3 points lost to the factorisation, and it applies
equally to every placement we compare.

\subsection{Crossbar Simulator}
We simulate the chip cell by cell in PyTorch. Each factorised layer becomes two crossbars per tier, one
holding $\sigma V^{\top}$ and one holding $U$, and every weight is stored in individual cells that can be
perturbed.

\subsubsection{Mapping a weight to cells}
A weight matrix $W$ is quantised to signed 8-bit integers $q = \mathrm{round}(W/s_i)$ with one scale
$s_i=\max_j|W_{ij}|/127$ per output row $i$. Because a conductance cannot be negative, each integer is
stored as a differential pair, $q=q^{+}-q^{-}$ with $q^{\pm}\ge0$. Each magnitude is split into base-$L$
digits, $L=2^{b}$, where $b$ is the number of bits per cell:
\begin{equation}
d_s = \left\lfloor m/L^{s}\right\rfloor \bmod L, \qquad m\in\{q^{+},q^{-}\}.
\end{equation}
An SLC cell ($b=1$, $L=2$) needs 7 slices and an MLC cell ($b=2$, $L=4$) needs 4. Digit $d$ is programmed as
the conductance
\begin{equation}
g = g_{\min} + d\,\Delta, \qquad \Delta=\frac{1-g_{\min}}{L-1},
\end{equation}
with conductances normalised so that $G_{\max}=1$ and $g_{\min}=1/100$ (an on/off ratio of 100). The
levels of an MLC cell are therefore three times closer together than those of an SLC cell. At read-out the
digit is recovered as $\hat d_s=(g^{+}_s-g^{-}_s)/\Delta$, and the slices are combined by shift-and-add,
$\hat W_{i\cdot}=s_i\sum_s L^{s}\hat d_s$.

\subsubsection{Error models}
Each error is applied to the programmed conductances, and the same random draw is used by every placement
that is evaluated on the same chip.
\begin{itemize}
\item \textbf{Programming noise.} Every cell receives $g \leftarrow \max\!\big(0,\;g+\varsigma\,\mathcal{N}(0,1)\big)$,
where $\varsigma$ is a fraction of $G_{\max}$. Since $\Delta$ is smaller for MLC, the same $\varsigma$
corresponds to a larger digit error in MLC. In some experiments $\varsigma$ is scaled by a factor that
depends on the stored level.
\item \textbf{Stuck-at faults.} A cell is stuck with probability $p$. A fraction $a=0.5$ of stuck cells is
stuck at $G_{\max}$ (SA1) and the rest at $g_{\min}$ (SA0), whatever value was programmed. Faults are drawn
independently per cell by default. We also model clusters of $3\times3$ cells and failed bit lines.
\item \textbf{Read noise.} Gaussian noise of standard deviation proportional to $G_{\max}$ is added fresh at
every read, in contrast to programming noise, which is fixed once the chip is programmed.
\item \textbf{IR drop.} Wire resistance reduces each conductance to $g/(1+r\,g\,n)$, where $n$ counts the wire
segments between the cell and its driver and sense amplifier, and $r$ is the segment resistance relative to
$R_{\mathrm{on}}$.
\end{itemize}

\subsubsection{Read-out and chip definition}
The input is quantised to 8 bits with one scale per tensor, as described above. Arrays are $128\times128$
tiles. By default the read-out is ideal. In the ADC experiments each slice of each tile is digitised
separately by an ADC whose range is fixed once on a noise-free chip over 64 calibration sentences, with
inputs applied bit-serially in sign-magnitude form. A \emph{chip} is one random draw of all programming
noise and faults, identified by a seed. All placements are evaluated on the same set of chips, so the
difference between two placements on a given chip does not come from different random draws. All
simulations run in FP32.

\subsection{Placement Policies}
A placement decides which SVD components are stored in the SLC tier and which in the MLC tier. The
factorised model is the same for every placement. Given an SLC share $f$, each layer with $k$ components
stores the $\lceil fk\rceil$ highest-scoring components in SLC and the remaining ones in MLC. Selection is
done within each layer, so every layer keeps the same SLC share. We compare six policies.
\begin{itemize}
\item \textbf{All-SLC} and \textbf{all-MLC} store every component in one tier. They are the most reliable
and the densest chip, and they bound the other policies.
\item \textbf{Random} selects a random set of components. It is a control that shows what a placement
without information achieves. The set is fixed, and only the cell errors differ between chips.
\item \textbf{Magnitude} scores component $r$ by its singular value $\sigma_r$. The largest singular values
carry most of the weight matrix.
\item \textbf{Gradient} follows HyFlexPIM \cite{song2025hyflexpim}. The score is
\begin{equation}
I_r=\sum_{b}\left|\frac{\partial L_b}{\partial\sigma_r}\right|,
\end{equation}
where $L_b$ is the loss on batch $b$ of 32 training sentences (pairs for MNLI) with dropout enabled. The absolute value
is taken per batch, because on a trained model the gradients of different batches point in different
directions and largely cancel in a single sum.
\item \textbf{Fisher} scores a component by the expected rise in loss if it is stored in MLC:
\begin{equation}
D_r=\sum_{e\in\mathcal{E}_r}F_e\,\mathrm{Var}(\epsilon_e),
\end{equation}
where $\mathcal{E}_r$ contains every stored entry of the component (row $r$ of $\sigma V^{\top}$ and column
$r$ of $U$) and $F_e$ is the empirical Fisher information of entry $e$, estimated from squared
per-sentence gradients over 2,048 training sentences. A fixed amount of cell noise becomes weight noise in
proportion to the 8-bit scale of the row that holds the entry, so $\mathrm{Var}(\epsilon_e)$ is proportional
to the squared scale of that row. Unlike the gradient, this score covers all entries disturbed by chip noise,
not only $\sigma_r$.
\end{itemize}
The calibration sentences are the same for the gradient and Fisher scores. Both are computed once per model
and are not affected by the chip noise. Table~\ref{tab:policies} summarises the policies.

\begin{table}[t]
\caption{Component placement policies. All policies store the same factorised model and differ only in the
score that selects the SLC components.}
\label{tab:policies}
\centering
\small
\begin{tabular}{@{}lp{1.45in}c@{}}
\toprule
Policy & Score of component $r$ & Uses loss \\
\midrule
All-SLC    & none (everything in SLC)                    & no  \\
All-MLC    & none (everything in MLC)                    & no  \\
Random     & random draw                                 & no  \\
Magnitude  & singular value $\sigma_r$                   & no  \\
Gradient   & $\sum_b|\partial L_b/\partial\sigma_r|$     & yes \\
Fisher     & $\sum_e F_e\,\mathrm{Var}(\epsilon_e)$ over all stored entries & yes \\
\bottomrule
\end{tabular}
\end{table}

\subsection{Fault Analysis}
HyFlexPIM models only programming noise. Real ReRAM arrays also contain stuck-at cells, which come from
forming failures and endurance wear-out. A stuck cell ignores the programmed value, so its error is not
small and is not centred on the target. We therefore study faults separately from noise in two
experiments. Neither adds programming noise, so every accuracy loss
comes from the stuck cells.
\begin{itemize}
\item \textbf{Fault rate.} The stuck-cell probability $p$ is swept over 0, 0.1, 0.2, 0.5, 1, 2 and 5\%, with
faults drawn independently per cell in both tiers. All six placements are evaluated at an SLC share of
20\% on 20 chips per point.
\item \textbf{Fault pattern.} At $p=1\%$, the stuck cells are placed in one of three ways, each with the
same expected number of stuck cells:
\emph{scattered}, every cell fails independently, as in the fault-rate sweep;
\emph{clustered}, blocks of $3\times3$ neighbouring cells fail together, as from a local defect; and
\emph{bit line}, all 128 cells of one column of a tile fail together, as from a broken wire or sense
amplifier.
In all cases half of the stuck cells are stuck at the highest conductance and half at the lowest.
All placements except random are evaluated at 20\% SLC on 10 chips per pattern.
\end{itemize}

\section{Results}
\subsection{Criticality-Aware Component Placement}
We first compare the placement policies as the share of components stored in SLC grows. Tables~\ref{tab:slc_distil} and \ref{tab:slc_bert} give the SST-2 accuracy of DistilBERT and BERT-Base at programming noise $\varsigma$ of 3\%, 5\% and 7.5\% of $G_{\max}$, averaged over 20 chips. All policies are evaluated on the same simulated chips. For reference, each block also lists the accuracy of the all-MLC and all-SLC chips at that noise level.

\begin{table}[t]
\caption{\textbf{DistilBERT}: SST-2 accuracy (\%) vs.\ SLC share (mean of 20 chips; bold = best).}
\label{tab:slc_distil}
\centering
\small
\begin{tabular}{@{}lccccc@{}}
\toprule
Policy & 5\% & 10\%  & 20\%  & 30\%  & 50\% \\
\midrule
\multicolumn{6}{@{}l}{\textit{Noise $\varsigma=3\%$}\quad all-MLC 86.0, all-SLC 88.7} \\
Random & 86.4 & 86.8 & 87.5 & 88.1 & 88.4 \\
Gradient & 88.4 & 88.5 & 88.6 & 88.8 & \textbf{88.7} \\
Magnitude & 88.4 & 88.6 & \textbf{88.8} & \textbf{88.9} & \textbf{88.7} \\
Fisher & \textbf{88.5} & \textbf{88.7} & 88.7 & 88.8 & 88.6 \\
\midrule
\multicolumn{6}{@{}l}{\textit{Noise $\varsigma=5\%$}\quad all-MLC 74.8, all-SLC 88.2} \\
Random & 75.4 & 77.4 & 81.9 & 82.3 & 85.7 \\
Gradient & 87.3 & 87.3 & 87.7 & 88.1 & 87.8 \\
Magnitude & \textbf{87.4} & \textbf{88.0} & 88.1 & \textbf{88.8} & \textbf{88.4} \\
Fisher & 87.3 & 87.8 & \textbf{88.4} & 88.1 & 88.3 \\
\midrule
\multicolumn{6}{@{}l}{\textit{Noise $\varsigma=7.5\%$}\quad all-MLC 56.9, all-SLC 87.4} \\
Random & 56.4 & 58.2 & 63.1 & 64.7 & 73.7 \\
Gradient & 78.0 & 80.7 & 83.1 & 84.8 & 86.0 \\
Magnitude & 80.2 & \textbf{84.0} & 85.2 & \textbf{86.6} & 86.9 \\
Fisher & \textbf{80.8} & 83.8 & \textbf{85.4} & 86.0 & \textbf{87.0} \\
\bottomrule
\end{tabular}
\end{table}

\begin{table}[t]
\caption{\textbf{BERT-Base}: SST-2 accuracy (\%) vs.\ SLC share (mean of 20 chips; bold = best).}
\label{tab:slc_bert}
\centering
\small
\begin{tabular}{@{}lccccc@{}}
\toprule
Policy & 5\% & 10\%  & 20\%  & 30\%  & 50\% \\
\midrule
\multicolumn{6}{@{}l}{\textit{Noise $\varsigma=3\%$}\quad all-MLC 82.1, all-SLC 89.2} \\
Random & 87.0 & 86.9 & 87.6 & 87.3 & 89.0 \\
Gradient & 88.7 & \textbf{90.0} & \textbf{90.0} & \textbf{89.9} & 89.6 \\
Magnitude & \textbf{89.7} & 89.9 & 89.8 & 89.8 & 89.7 \\
Fisher & 89.6 & 89.8 & 89.8 & 89.8 & \textbf{89.8} \\
\midrule
\multicolumn{6}{@{}l}{\textit{Noise $\varsigma=5\%$}\quad all-MLC 66.8, all-SLC 88.0} \\
Random & 72.5 & 73.6 & 76.2 & 77.8 & 83.4 \\
Gradient & 84.2 & 88.2 & 88.4 & 88.4 & 88.1 \\
Magnitude & \textbf{88.0} & \textbf{88.6} & 88.7 & \textbf{88.9} & \textbf{88.9} \\
Fisher & 87.9 & \textbf{88.6} & \textbf{88.8} & \textbf{88.9} & \textbf{88.9} \\
\midrule
\multicolumn{6}{@{}l}{\textit{Noise $\varsigma=7.5\%$}\quad all-MLC 53.0, all-SLC 85.7} \\
Random & 53.6 & 54.2 & 55.2 & 57.5 & 67.3 \\
Gradient & 67.5 & 79.0 & 81.2 & 82.4 & 83.2 \\
Magnitude & 79.4 & 81.8 & 84.1 & \textbf{85.5} & 85.6 \\
Fisher & \textbf{80.0} & \textbf{82.2} & \textbf{84.2} & 85.1 & \textbf{85.8} \\
\bottomrule
\end{tabular}
\end{table}

Random placement recovers little of the gap between all-MLC and all-SLC, even at 50\% SLC, and its accuracy
varies strongly between chips. The three informed policies recover most of it with far less SLC. At 3\%
noise the gap is small (2.7 points on DistilBERT, 7.1 on BERT-Base), and every informed policy reaches the
all-SLC accuracy from 5--10\% SLC; the policies differ by less than the chip-to-chip spread. At 5\%
noise, magnitude and Fisher placement are within about one point of the all-SLC chip from 5\% SLC (10\% for
the gradient on BERT-Base). At 7.5\% noise, 20\% SLC leaves magnitude and Fisher about 2 points below
all-SLC on DistilBERT and 1.5 points below on BERT-Base. Magnitude and Fisher differ by less than the
chip-to-chip spread, so the data does not separate them. The gradient policy is the weakest of the three.
It trails magnitude by 2 points at 20\% SLC on DistilBERT and 3 points on BERT-Base, and by 12 points at 5\%
SLC on BERT-Base at 7.5\% noise.

\subsection{Fault Analysis}
Table~\ref{tab:fault} gives the accuracy against the share of stuck cells, and Table~\ref{tab:pattern} the
accuracy for the three fault patterns at $p=1\%$. At 5\% stuck cells every placement is at chance (about
50\%), so that column is omitted.

\begin{table}[t]
\caption{Accuracy (\%) vs.\ share of stuck cells, no programming noise (20\% SLC, mean of 20 chips).}
\label{tab:fault}
\centering
\small
\begin{tabular}{@{}lcccccc@{}}
\toprule
Policy & 0 & 0.1\% & 0.2\% & 0.5\% & 1\% & 2\% \\
\midrule
\multicolumn{7}{@{}l}{\textbf{DistilBERT}} \\
All-SLC   & 89.0 & 88.2 & 87.8 & 85.2 & 74.6 & 58.5 \\
All-MLC   & 89.1 & 65.5 & 58.2 & 52.3 & 51.1 & 50.2 \\
Random    & 88.6 & 76.9 & 63.3 & 51.3 & 51.0 & 49.4 \\
Gradient  & 89.8 & 88.1 & 86.7 & 74.5 & 55.4 & 51.0 \\
Magnitude & 89.1 & 88.3 & 87.4 & 80.3 & 62.1 & 51.4 \\
Fisher    & 89.0 & 88.3 & 87.2 & 80.3 & 62.1 & 51.3 \\
\midrule
\multicolumn{7}{@{}l}{\textbf{BERT-Base}} \\
All-SLC   & 89.9 & 88.3 & 85.6 & 79.0 & 62.7 & 51.4 \\
All-MLC   & 90.0 & 58.6 & 52.6 & 50.8 & 51.2 & 49.6 \\
Random    & 90.0 & 67.9 & 55.8 & 50.7 & 50.6 & 50.2 \\
Gradient  & 89.8 & 89.4 & 87.5 & 71.0 & 52.8 & 50.2 \\
Magnitude & 89.7 & 89.6 & 88.6 & 76.8 & 57.9 & 51.0 \\
Fisher    & 89.1 & 89.7 & 88.4 & 75.6 & 55.5 & 49.5 \\
\bottomrule
\end{tabular}
\end{table}

\begin{table}[t]
\caption{Accuracy (\%) at 1\% stuck cells for three fault patterns (20\% SLC, mean of 10 chips).}
\label{tab:pattern}
\centering
\small
\begin{tabular}{@{}lcccccc@{}}
\toprule
 & \multicolumn{3}{c}{\textbf{DistilBERT}} & \multicolumn{3}{c}{\textbf{BERT-Base}} \\
\cmidrule(lr){2-4}\cmidrule(l){5-7}
Policy & Scat. & Clust. & Line & Scat. & Clust. & Line \\
\midrule
All-SLC   & 74.2 & 80.3 & 71.2 & 61.3 & 66.7 & 57.0 \\
All-MLC   & 51.2 & 51.2 & 50.8 & 50.7 & 50.8 & 50.4 \\
Gradient  & 55.0 & 56.4 & 52.2 & 52.6 & 51.5 & 51.0 \\
Magnitude & 62.5 & 56.1 & 54.7 & 56.4 & 54.0 & 50.4 \\
Fisher    & 60.0 & 57.3 & 56.3 & 55.6 & 53.1 & 52.2 \\
\bottomrule
\end{tabular}
\end{table}

MLC is far more fragile to faults than to noise. Only 0.1\% stuck cells reduce the all-MLC accuracy by 24
points on DistilBERT and 31 points on BERT-Base, whereas programming noise of 2\% on every cell costs it
about 1 and 3 points. Random placement fails almost as early, because it leaves important components in MLC.
With magnitude or Fisher placement, the hybrid chip stays within 2 points of the fault-free accuracy up to
0.2\% stuck cells, and on BERT-Base it is even more accurate than the all-SLC chip there, since SLC needs
almost twice as many cells per weight. From 0.5\% the hybrid chip loses 9--14 points, the gradient
placement 15--19 points, and at 1\% even the all-SLC chip fails. Clustered faults hurt the all-SLC chip less
than scattered ones, since they damage fewer weights, while failed bit lines hurt most. At 1\% stuck cells
every hybrid placement is close to chance for all three patterns.

\subsection{Fault-Aware Mapping}
Placement cannot protect against faults, because a stuck cell damages whichever component is stored on it.
The SVD structure, however, offers a fix at no hardware cost. Component $r$ occupies one row of the
$\sigma V^{\top}$ crossbar and the matching input row of the $U$ crossbar, and the components of a layer can
be stored in any order without changing its output. After a chip has been tested and its stuck cells are
known, we choose which component goes into which physical slot. The cost of storing component $r$ in slot
$s$ is its expected rise in loss,
\begin{equation}
C_{rs}=\sum_{e\in\mathcal{E}_r}F_e\,\mathbb{E}\big[\epsilon_{e,s}^2\big],
\end{equation}
where $\epsilon_{e,s}$ is the weight error of entry $e$ in slot $s$: the exact error caused by the slot's
stuck cells plus the programming noise of the slot's tier. The Hungarian algorithm finds the assignment of
lowest total cost. We test three mapped placements against the unmapped ones on the same 10 chips:
\emph{all-SLC + remap}, \emph{magnitude + remap}, which keeps the 20\% magnitude SLC set and reorders the
components within each tier, and \emph{C1}, which chooses the tier and the slot of every component in a
single assignment. Table~\ref{tab:remap} gives the results.

\begin{table}[t]
\caption{Fault-aware mapping: accuracy (\%) vs.\ share of stuck cells and fault pattern, no programming
noise (mean of 10 chips).}
\label{tab:remap}
\centering
\small
\begin{tabular}{@{}lcccccc@{}}
\toprule
 & \multicolumn{4}{c}{Scattered} & \multicolumn{2}{c}{1\%} \\
\cmidrule(lr){2-5}\cmidrule(l){6-7}
Policy & 0.2\% & 0.5\% & 1\% & 2\% & Clust. & Line \\
\midrule
\multicolumn{7}{@{}l}{\textbf{DistilBERT}} \\
All-SLC           & 88.1 & 85.4 & 73.7 & 57.9 & 79.6 & 71.1 \\
All-SLC + remap   & 88.0 & \textbf{87.6} & \textbf{84.8} & \textbf{62.3} & \textbf{85.6} & \textbf{81.5} \\
Magnitude         & 87.1 & 77.9 & 62.4 & 50.2 & 55.9 & 54.0 \\
Magnitude + remap & \textbf{88.4} & 85.1 & 67.0 & 53.3 & 68.5 & 68.0 \\
C1                & 88.3 & 85.3 & 67.7 & 54.7 & 75.5 & 62.7 \\
\midrule
\multicolumn{7}{@{}l}{\textbf{BERT-Base}} \\
All-SLC           & 85.4 & 80.1 & 60.9 & 51.1 & 66.5 & 57.2 \\
All-SLC + remap   & 89.6 & 83.7 & \textbf{75.8} & \textbf{55.9} & \textbf{81.3} & \textbf{77.5} \\
Magnitude         & 88.5 & 76.6 & 55.9 & 50.1 & 53.7 & 50.1 \\
Magnitude + remap & \textbf{90.1} & \textbf{85.6} & 65.1 & 52.2 & 64.3 & 56.9 \\
C1                & 90.0 & 84.3 & 68.3 & 49.4 & 61.5 & 56.8 \\
\bottomrule
\end{tabular}
\end{table}

Remapping improves every placement it is applied to. On the same chips, magnitude + remap beats magnitude at
every fault rate, pattern and model. At 0.5\% stuck cells it raises the accuracy from 77.9 to 85.1\% on
DistilBERT and from 76.6 to 85.6\% on BERT-Base, and for clustered faults it gains 10--13 points. The
all-SLC chip gains most at higher fault rates: at 1\% stuck cells it rises from 73.7 to 84.8\% on DistilBERT
and from 60.9 to 75.8\% on BERT-Base, and for failed bit lines by 10 and 20 points. Up to 0.2\% stuck cells,
the remapped hybrid chip matches the remapped all-SLC chip with about 34\% fewer cells; from 0.5\% the
remapped all-SLC chip is clearly more robust. Choosing the tier per chip (C1) gives no reliable gain over
magnitude + remap: it wins in some settings and loses in others. The gains hold when 5\% programming noise
is added: at 0.5\% stuck cells remapping raises magnitude placement from 74.9 to 83.0\% on DistilBERT and
from 70.3 to 80.4\% on BERT-Base.

\subsection{IR Drop and ADC Resolution}
We next test two non-idealities of the array periphery: wire resistance (IR drop) and a fixed ADC. In both
experiments there is no programming noise and no fault, so the result is deterministic and each setting is
evaluated on a single chip, with hybrid placements at 20\% SLC. Differences below about one point are
within the spread between placements of an ideal chip and are not meaningful.

\textbf{IR drop.} The wire-to-cell resistance ratio $r$ is swept from $0$ to $3\times10^{-3}$ with no
compensation. Table~\ref{tab:irdrop} gives the results.

\begin{table}[t]
\caption{IR drop: accuracy (\%) vs.\ wire resistance ratio $r$ (20\% SLC, no programming noise).}
\label{tab:irdrop}
\centering
\small
\begin{tabular}{@{}lccccc@{}}
\toprule
Policy & 0 & $10^{-4}$ & $3{\cdot}10^{-4}$ & $10^{-3}$ & $3{\cdot}10^{-3}$ \\
\midrule
\multicolumn{6}{@{}l}{\textbf{DistilBERT}} \\
All-MLC   & 89.0 & 88.6 & 89.1 & \textbf{89.9} & \textbf{77.9} \\
All-SLC   & 88.5 & 89.9 & 89.4 & 85.3 & 70.1 \\
Magnitude & 88.8 & 89.2 & 89.7 & 86.6 & 73.4 \\
Gradient  & 88.6 & 89.6 & 89.6 & 86.9 & 74.7 \\
Fisher    & 89.1 & 89.2 & 89.3 & 86.4 & 72.6 \\
\midrule
\multicolumn{6}{@{}l}{\textbf{BERT-Base}} \\
All-MLC   & 90.4 & 89.3 & 89.8 & \textbf{90.6} & \textbf{79.1} \\
All-SLC   & 89.7 & 90.3 & 88.9 & 87.0 & 58.5 \\
Magnitude & 90.6 & 90.1 & 90.4 & 88.1 & 76.8 \\
Gradient  & 90.3 & 90.5 & 90.5 & 88.6 & 76.4 \\
Fisher    & 89.6 & 91.1 & 91.1 & 88.3 & 72.8 \\
\bottomrule
\end{tabular}
\end{table}

Up to $r=3\times10^{-4}$ no placement loses accuracy. At $r=10^{-3}$ all-MLC is unaffected while all-SLC
loses 3 points, and at $r=3\times10^{-3}$ all-MLC keeps 78--79\% while all-SLC falls to 70\% on DistilBERT
and 59\% on BERT-Base. IR drop therefore reverses the usual order: SLC is the more sensitive tier. The
attenuation $r\,g\,n$ grows with the conductance, and SLC stores every non-zero digit at $G_{\max}$. The
hybrid placements lie between the two, with no consistent difference between policies.

\textbf{ADC resolution.} Inputs are applied bit-serially and each slice of each tile is read by an ADC of
3 to 8 bits with a fixed range. Table~\ref{tab:adc} gives the results down to 5 bits; at 4 and 3 bits every
placement is near chance.

\begin{table}[t]
\caption{ADC: accuracy (\%) vs.\ ADC resolution in bits (20\% SLC, no programming noise).}
\label{tab:adc}
\centering
\small
\begin{tabular}{@{}lccccc@{}}
\toprule
Policy & Ideal & 8 & 7 & 6 & 5 \\
\midrule
\multicolumn{6}{@{}l}{\textbf{DistilBERT}} \\
All-MLC   & 89.0 & 88.3 & 87.8 & 81.7 & 56.5 \\
All-SLC   & 88.5 & 89.2 & 88.4 & 88.0 & 71.0 \\
Magnitude & 88.8 & 88.6 & 89.0 & 89.7 & 70.6 \\
Gradient  & 88.6 & 88.4 & 88.5 & 88.4 & 77.3 \\
Fisher    & 89.1 & 88.0 & 88.2 & 88.8 & 70.9 \\
\midrule
\multicolumn{6}{@{}l}{\textbf{BERT-Base}} \\
All-MLC   & 90.4 & 89.6 & 89.2 & 83.0 & 51.1 \\
All-SLC   & 89.7 & 89.9 & 90.0 & 89.3 & 51.1 \\
Magnitude & 90.6 & 90.4 & 89.3 & 90.0 & 51.1 \\
Gradient  & 90.3 & 90.1 & 90.1 & 89.6 & 53.0 \\
Fisher    & 89.6 & 90.1 & 89.7 & 90.5 & 51.6 \\
\bottomrule
\end{tabular}
\end{table}

A 7-bit ADC costs no accuracy for any placement. At 6 bits all-MLC loses 7 points on both models, while
all-SLC and every hybrid placement remain within one point of the ideal read-out. At 5 bits every
placement fails on BERT-Base, and on DistilBERT all keep only 71--77\%. \textbf{MLC needs at least one
more ADC bit than SLC, so the hybrid chip can run with a 6-bit ADC.} With
5\% programming noise added (DistilBERT, 3 chips), the gap widens: with a 6-bit ADC all-MLC drops from
77.8\% (ideal read-out) to 56.0\%, while all-SLC and the hybrid placements stay at 87--89\%.

\subsection{Second Task: MNLI}
We repeat three experiments on MNLI: placement at $\varsigma=5\%$, the fault-rate sweep, and fault-aware
mapping. IR drop, ADC resolution, fault patterns and the other noise levels are not repeated. MNLI has three
classes, so chance is about 33\%.

\begin{table}[t]
\caption{\textbf{MNLI}: accuracy (\%) in FP32 and INT8, with no hardware error (2,000-pair subset; hybrid
placements at 20\% SLC).}
\label{tab:mnli_int8}
\centering
\small
\begin{tabular}{@{}llcc@{}}
\toprule
 & & DistilBERT & BERT-Base \\
\midrule
FP32 & Original          & 82.1 & 83.7 \\
     & SVD + fine-tune   & 80.9 & 83.2 \\
\midrule
INT8 & All-SLC           & 76.9 & 72.2 \\
     & All-MLC           & 76.8 & 72.4 \\
     & Random            & 77.7 & 73.0 \\
     & Gradient          & \textbf{79.6} & 78.2 \\
     & Magnitude         & 77.3 & 79.1 \\
     & Fisher            & 78.3 & \textbf{79.3} \\
\bottomrule
\end{tabular}
\end{table}

Unlike SST-2, MNLI is sensitive to the 8-bit quantisation (Table~\ref{tab:mnli_int8}). With no noise and no
faults, the all-SLC chip loses 4 points against the FP32 SVD model on DistilBERT (80.9
to 76.9\%) and 11 points on BERT-Base (83.2 to 72.2\%). This loss also depends on the placement. Each tier
has its own 8-bit scales, so a placement that separates large from small components quantises the small
ones more finely. On BERT-Base the noise-free magnitude, gradient and Fisher chips reach 78--79\%, while the
all-SLC, all-MLC and random chips stay at 72--73\%. Part of the gap between the hybrid and the all-SLC chip
on BERT-Base therefore comes from quantisation and not from cell errors.

\begin{table}[t]
\caption{\textbf{MNLI}: accuracy (\%) vs.\ SLC share at $\varsigma=5\%$ (mean of 20 chips; bold = best).}
\label{tab:mnli_slc}
\centering
\small
\begin{tabular}{@{}lccccc@{}}
\toprule
Policy & 5\% & 10\%  & 20\%  & 30\%  & 50\% \\
\midrule
\multicolumn{6}{@{}l}{\textbf{DistilBERT}\quad all-MLC 36.8, all-SLC 72.0} \\
Random    & 38.3 & 39.3 & 44.1 & 48.1 & 55.3 \\
Gradient  & 66.6 & 69.0 & 71.3 & 72.1 & \textbf{73.6} \\
Magnitude & 68.4 & \textbf{70.0} & \textbf{71.5} & \textbf{72.6} & 72.5 \\
Fisher    & \textbf{68.5} & \textbf{70.0} & 71.1 & 72.1 & 72.3 \\
\midrule
\multicolumn{6}{@{}l}{\textbf{BERT-Base}\quad all-MLC 33.5, all-SLC 61.6} \\
Random    & 34.1 & 33.9 & 35.9 & 38.9 & 44.9 \\
Gradient  & 59.6 & 63.5 & 63.3 & 66.5 & 64.2 \\
Magnitude & 64.0 & \textbf{66.0} & \textbf{68.0} & \textbf{69.2} & \textbf{66.7} \\
Fisher    & \textbf{64.2} & 65.8 & 67.8 & 68.8 & 66.6 \\
\bottomrule
\end{tabular}
\end{table}

Table~\ref{tab:mnli_slc} gives the placement results. All-MLC is at chance on both models, and random
placement recovers little even at 50\% SLC. The informed policies again recover most of the gap with little
SLC. On DistilBERT, magnitude and Fisher placement are within 1 point of the all-SLC chip from 20\% SLC. On
BERT-Base, magnitude and Fisher beat the all-SLC chip from 5\% SLC, by up to 7.6 points at 30\%. From 30\% to
50\% SLC, the accuracy of every informed policy on BERT-Base falls by about 2.5 points. Magnitude and Fisher
again differ by less than the chip-to-chip spread. The gradient policy is again the weakest of the three.
It trails magnitude by 4.7 points at 20\% SLC on BERT-Base and by 4.4 points at 5\% SLC.

\begin{table}[t]
\caption{\textbf{MNLI}: agreement between component scores, mean over layers (rank correlation $\rho$ and
overlap of the top 20\%).}
\label{tab:mnli_scores}
\centering
\small
\begin{tabular}{@{}lcccc@{}}
\toprule
 & \multicolumn{2}{c}{\textbf{DistilBERT}} & \multicolumn{2}{c}{\textbf{BERT-Base}} \\
\cmidrule(lr){2-3}\cmidrule(l){4-5}
Scores & $\rho$ & Top 20\% & $\rho$ & Top 20\% \\
\midrule
Fisher vs.\ magnitude   & 0.97 & 89\% & 0.96 & 88\% \\
Gradient vs.\ magnitude & 0.57 & 67\% & 0.56 & 63\% \\
Gradient vs.\ Fisher    & 0.66 & 71\% & 0.66 & 69\% \\
\bottomrule
\end{tabular}
\end{table}

Table~\ref{tab:mnli_scores} explains these results. The Fisher score ranks the components almost as the
singular values do, and it shares 88--89\% of its top 20\% with magnitude. The gradient score agrees much
less with both, and about a third of its top 20\% is different from the magnitude selection.

\begin{table}[t]
\caption{\textbf{MNLI}: accuracy (\%) vs.\ share of stuck cells, no programming noise (20\% SLC, mean of 20
chips).}
\label{tab:mnli_fault}
\centering
\small
\begin{tabular}{@{}lccccc@{}}
\toprule
Policy & 0 & 0.1\% & 0.2\% & 0.5\% & 1\% \\
\midrule
\multicolumn{6}{@{}l}{\textbf{DistilBERT}} \\
All-SLC   & 76.9 & 74.4 & 71.3 & 60.4 & 42.1 \\
All-MLC   & 76.8 & 44.5 & 34.0 & 33.6 & 33.7 \\
Random    & 77.7 & 51.6 & 35.3 & 32.9 & 33.2 \\
Gradient  & 79.6 & 76.0 & 71.4 & 44.5 & 33.9 \\
Magnitude & 77.3 & 74.8 & 71.3 & 49.7 & 34.0 \\
Fisher    & 78.3 & 74.9 & 70.7 & 50.3 & 34.0 \\
\midrule
\multicolumn{6}{@{}l}{\textbf{BERT-Base}} \\
All-SLC   & 72.2 & 67.2 & 61.9 & 44.4 & 34.5 \\
All-MLC   & 72.4 & 37.2 & 33.2 & 32.8 & 31.8 \\
Random    & 73.0 & 42.8 & 34.3 & 33.3 & 32.6 \\
Gradient  & 78.2 & 73.0 & 65.2 & 36.9 & 32.0 \\
Magnitude & 79.1 & 74.6 & 67.8 & 40.4 & 32.8 \\
Fisher    & 79.3 & 73.9 & 67.2 & 39.8 & 32.8 \\
\bottomrule
\end{tabular}
\end{table}

Table~\ref{tab:mnli_fault} gives the fault-rate sweep. At 2\% and 5\% stuck cells every placement is at
chance, so these columns are omitted. MNLI is more fragile to stuck cells than SST-2. All-MLC and random
placement fail at 0.1\% stuck cells. Magnitude and Fisher placement lose 2--5 points at 0.1\% and 6--11
points at 0.2\%, and on BERT-Base they stay above the all-SLC chip (67.8 against 61.9\% at 0.2\%). From
0.5\% every placement loses 16 points or more, with the gradient placement lowest among the hybrid ones, and
at 1\% every placement is near chance, including all-SLC.

\begin{table}[t]
\caption{\textbf{MNLI}: fault-aware mapping, accuracy (\%) vs.\ share of stuck cells, no programming noise
(mean of 10 chips).}
\label{tab:mnli_remap}
\centering
\small
\begin{tabular}{@{}lcccccc@{}}
\toprule
 & \multicolumn{3}{c}{\textbf{DistilBERT}} & \multicolumn{3}{c}{\textbf{BERT-Base}} \\
\cmidrule(lr){2-4}\cmidrule(l){5-7}
Policy & 0.2\% & 0.5\% & 1\% & 0.2\% & 0.5\% & 1\% \\
\midrule
All-SLC           & 71.4 & 58.4 & 41.4 & 61.5 & 46.7 & 35.7 \\
All-SLC + remap   & 76.2 & \textbf{72.2} & \textbf{58.5} & 69.9 & \textbf{61.9} & \textbf{46.2} \\
Magnitude         & 69.9 & 49.1 & 33.4 & 68.1 & 40.0 & 33.1 \\
Magnitude + remap & \textbf{76.6} & 67.9 & 38.8 & 75.3 & 60.0 & 36.5 \\
C1                & \textbf{76.6} & 67.0 & 39.4 & \textbf{75.5} & 60.4 & 35.5 \\
\bottomrule
\end{tabular}
\end{table}

Table~\ref{tab:mnli_remap} gives the mapping results. At 2\% stuck cells every placement is at chance
(32--37\%), so that column is omitted. Remapping again improves every placement. At 0.2\% stuck cells,
magnitude + remap raises the accuracy from 69.9 to 76.6\% on DistilBERT and from 68.1 to 75.3\% on
BERT-Base, and it cuts the chip-to-chip standard deviation from 4.5 to 0.4 points and from 3.4 to 0.9
points. From 0.5\% the remapped all-SLC chip is the most robust: at 1\% stuck cells it gains 17 points over
all-SLC on DistilBERT and 11 points on BERT-Base. C1 again gives no reliable gain over magnitude + remap.

% ---------------------------------------------------------------
% Remaining results, discussion and conclusion to be written.
% ---------------------------------------------------------------

\begin{thebibliography}{00}
\bibitem{song2025hyflexpim} C.~E. Song, P.~Bhatnagar, Z.~Xia, N.~S. Kim, T.~S. Rosing, and M.~Kang,
``Hybrid SLC-MLC RRAM mixed-signal processing-in-memory architecture for transformer acceleration via
gradient redistribution,'' in \emph{Proc. 52nd Annu. Int. Symp. Computer Architecture (ISCA)}, Tokyo,
Japan, Jun. 2025, pp. 1155--1170.
\bibitem{devlin2019bert} J.~Devlin, M.-W. Chang, K.~Lee, and K.~Toutanova, ``BERT: Pre-training of deep
bidirectional transformers for language understanding,'' in \emph{Proc. NAACL-HLT}, 2019.
\bibitem{sanh2019distilbert} V.~Sanh, L.~Debut, J.~Chaumond, and T.~Wolf, ``DistilBERT, a distilled version of
BERT: smaller, faster, cheaper and lighter,'' \emph{arXiv preprint arXiv:1910.01108}, 2019.
\bibitem{wang2019glue} A.~Wang, A.~Singh, J.~Michael, F.~Hill, O.~Levy, and S.~R. Bowman, ``GLUE: A
multi-task benchmark and analysis platform for natural language understanding,'' in \emph{Proc. ICLR}, 2019.
\end{thebibliography}

\end{document}
